\chapter{मुख्य आदर्श प्रांत पर मॉड्यूल}\label{ch:b3:modules}

क्षेत्र के स्थान पर किसी \emph{वलय} पर रैखिक बीजगणित: परिकल्पना
का यह छोटा-सा परिवर्तन बीजगणित की महान एकीकरण प्रमेयों में से एक
को जन्म देता है। $\Z$ पर \omterm{def:b3:modules:module}{मॉड्यूल} एक आबेली समूह है; $K[X]$ पर \omterm{def:b3:modules:module}{मॉड्यूल}
एक ऐसी सदिश समष्टि है जिसके साथ एक अंतःसमाकारिता जुड़ी है। अतः
किसी \omterm{def:b3:rings:pidufd}{PID} पर \omterm{def:b3:modules:free}{परिमिततः जनित} मॉड्यूलों की संरचना प्रमेय एक ही आघात
में समस्त \omterm{def:b3:modules:free}{परिमिततः जनित} आबेली समूहों \emph{और} समरूपता तक समस्त
अंतःसमाकारिताओं का वर्गीकरण कर देती है --- जॉर्डन अपचयन, जिसे
दूसरे वर्ष में नाजुक आगमनों से प्राप्त किया गया था, यहाँ एक
उपप्रमेय के रूप में झड़ पड़ता है, और उसके साथ उसका अधिक सूक्ष्म
सहोदर भी: \emph{परिमेय} विहित रूप, जो प्रत्येक क्षेत्र पर मान्य
है। परिकलन का यंत्र है \omterm{thm:b3:modules:smith}{स्मिथ प्रसामान्य रूप} --- आव्यूहों का एक
ऐसा अंकगणित जो यूक्लिड के योग्य है।

सर्वत्र $A$ एक क्रमविनिमेय वलय है, और शीघ्र ही एक \omterm{def:b3:rings:pidufd}{PID};
``\omterm{def:b3:modules:module}{मॉड्यूल}'' से अभिप्राय है $A$-मॉड्यूल।

\section{मॉड्यूल, मुक्त मॉड्यूल}

\begin{definition}\label{def:b3:modules:module}
कोई \emph{$A$-मॉड्यूल}\index{मॉड्यूल} एक आबेली समूह $(M, +)$ है जिस
पर अदिश गुणन $A \times M \to M$ सदिश-समष्टि स्वयंसिद्धों को संतुष्ट करता है:
$a(x + y) = ax + ay$, $(a + b)x = ax + bx$, $(ab)x = a(bx)$, $1x = x$। उपमॉड्यूल, भागफल $M/N$, समाकारिताएँ
($A$-रैखिक प्रतिचित्रण), प्रत्यक्ष योग $\bigoplus_i M_i$ और तुल्याकारिता
प्रमेय ठीक उसी प्रकार परिभाषित और सिद्ध होते हैं जैसे सदिश
समष्टियों और आबेली समूहों के लिए; विशेष रूप से किसी समाकारिता
$f$ के लिए $M/\ker f \cong
\operatorname{im} f$।
\end{definition}

\begin{example}\label{ex:b3:modules:examples}
तीन प्रेरक स्थितियाँ।
\begin{enumerate}
  \item $A = K$ कोई क्षेत्र: \omterm{def:b3:modules:module}{मॉड्यूल} सदिश समष्टियाँ हैं।
  \item $A = \Z$: \omterm{def:b3:modules:module}{मॉड्यूल} ठीक आबेली समूह हैं ($nx$ का $x + \dots + x$ होना
        अनिवार्य है), और उपमॉड्यूल उपसमूह हैं।
  \item $A = K[X]$: कोई \omterm{def:b3:modules:module}{मॉड्यूल} एक $K$-सदिश समष्टि $V$ है जिसके साथ
        $K$-रैखिक प्रतिचित्रण $u\colon x \mapsto X\cdot x$ जुड़ा है --- विलोमतः, $u \in \mathcal L(V)$
        वाला प्रत्येक युग्म $(V, u)$ $P \cdot x = P(u)(x)$ द्वारा $K[X]$-मॉड्यूल बन
        जाता है। उपमॉड्यूल ठीक $u$-स्थिर उपसमष्टियाँ हैं।
\end{enumerate}
$A$ का \omterm{def:b3:rings:ideal}{आदर्श} ठीक $A$ का उपमॉड्यूल है; \omterm{def:b3:rings:ideal}{भागफल वलय} $A/I$ एक
$A$-मॉड्यूल है। सदिश समष्टियों के विपरीत, मॉड्यूलों में
\emph{\omterm{def:b3:modules:torsion}{मरोड़}} हो सकती है: $\Z/6\Z$ में अवयव $\bar 3 \ne 0$ $2 \neq 0$ से नष्ट हो जाता
है।
\end{example}

\begin{definition}\label{def:b3:modules:free}
$M$ \emph{परिमिततः जनित}\index{परिमिततः जनित मॉड्यूल} कहलाता है
यदि किसी $x_i$ के लिए $M = Ax_1 + \dots + Ax_n$ हो। $M$ कोटि $n$ का
\emph{मुक्त}\index{मुक्त मॉड्यूल} कहलाता है यदि $M \cong A^n$, अर्थात्
यदि उसका कोई \emph{आधार} हो (ऐसा जनक कुल जो $A$-रैखिकतः स्वतंत्र
है)। प्रत्येक परिमिततः जनित $M$ किसी मुक्त \omterm{def:b3:modules:module}{मॉड्यूल} का भागफल
होता है: $(a_1, \dots, a_n) \mapsto \sum a_ix_i$ $A^n$ को $M$ पर आच्छादित करता है।
\end{definition}

\begin{proposition}[कोटि की अपरिवर्तिता]\label{prop:b3:modules:rank}
यदि $A \neq 0$ और $A^m \cong A^n$ हों, तो $m = n$।
\end{proposition}

\begin{proof}
$A$ का कोई \omterm{def:b3:rings:primemaximal}{महत्तम आदर्श} $\mathfrak m$ चुनिए (\cref{thm:b3:rings:krull}) और $k = A/\mathfrak m$ रखिए, जो एक
क्षेत्र है। कोई तुल्याकारिता $f \colon A^m
\to A^n$ रैखिकता से $\mathfrak m A^m$ को $\mathfrak m A^n$ में
भेजती है, अतः $k$-सदिश समष्टियों के रूप में भागफलों की
तुल्याकारिता
\[
A^m/\mathfrak m A^m \;\cong\; A^n/\mathfrak m A^n,
\qquad\text{अर्थात्}\qquad k^m \cong k^n
\]
प्रेरित करती है (भागफल $A^m/\mathfrak m A^m$ $\mathfrak m$ से नष्ट हो जाता है, अतः $A$
की क्रिया $k$ से होकर गुजरती है; मानक आधार के प्रतिबिंब एक
$k$-आधार बनाते हैं)। क्षेत्र $k$ पर विमा सिद्धांत से $m = n$।
\end{proof}

\begin{theorem}[मुक्त मॉड्यूलों के उपमॉड्यूल]\label{thm:b3:modules:submodule}
मान लीजिए $A$ एक \omterm{def:b3:rings:pidufd}{PID} है और $M \subseteq A^n$ एक उपमॉड्यूल। तब $M$ कोटि
$\leq n$ का \omterm{def:b3:modules:free}{मुक्त मॉड्यूल} है।
\end{theorem}

\begin{proof}
$n$ पर आगमन। $n = 1$ के लिए: $M$ एक \omterm{def:b3:rings:ideal}{आदर्श} है, अतः $M = (0)$ (कोटि
$0$ का \omterm{def:b3:modules:free}{मुक्त}) या $M = dA \cong A$ ($x \mapsto dx$ एकैकी है: प्रांत)। $n > 1$ के लिए:
मान लीजिए $\pi \colon A^n \to A$ अंतिम निर्देशांक है। तब $\pi(M)$ एक \omterm{def:b3:rings:ideal}{आदर्श} है, $(0)$
या $dA$। यदि $(0)$: $M \subseteq A^{n-1} \times \{0\}$ और आगमन लागू हो जाता है। अन्यथा $\pi(x_0) = d$
वाला $x_0 \in M$ चुनिए। प्रत्येक $x \in M$ अद्वितीय रूप से
\[
x = \underbrace{\frac{\pi(x)}{d}}_{\in A}\, x_0 +
\Bigl(x - \tfrac{\pi(x)}d x_0\Bigr),
\qquad x - \tfrac{\pi(x)}d x_0 \in M \cap \ker\pi
\]
लिखा जाता है ($\pi(x) \in dA$, अतः गुणांक $A$ में है)। इस प्रकार $M = Ax_0
\oplus (M \cap \ker \pi)$:
योग प्रत्यक्ष है, क्योंकि $\pi(ax_0) =
ad = 0$ से $a = 0$ अनिवार्य हो जाता है।
आगमन से $M \cap \ker\pi \subseteq
\ker \pi \cong A^{n-1}$ कोटि $\leq n - 1$ का \omterm{def:b3:modules:free}{मुक्त मॉड्यूल} है; उसमें $x_0$
(जो अभी देखे अनुसार $\ker\pi$ से स्वतंत्र है) जोड़ने पर गणनीयता
$\leq n$ वाला $M$ का आधार मिल जाता है।
\end{proof}

\begin{remark}\label{rem:b3:modules:presentation}
फलस्वरूप, किसी \omterm{def:b3:rings:pidufd}{PID} पर प्रत्येक \omterm{def:b3:modules:free}{परिमिततः जनित मॉड्यूल} $M$ की एक
\emph{परिमित प्रस्तुति} होती है: किसी आच्छादन $\varphi\colon A^n \to M$ की अष्टि
\omterm{def:b3:modules:free}{मुक्त} होती है और उसका आधार $c_1, \dots, c_k$ ($k \leq n$) होता है, तथा $M
\cong A^n / \operatorname{im}(C)$,
जहाँ $C \in M_{n,k}(A)$ वह आव्यूह है जिसके स्तंभ $c_j$ हैं। $M$ को समझने का
अर्थ है स्रोत और लक्ष्य में आधार-परिवर्तन तक \emph{$A$ पर एक
आव्यूह} को समझना --- यही अगले अनुभाग का विषय है।
\end{remark}

\section{स्मिथ प्रसामान्य रूप}

\begin{definition}\label{def:b3:modules:equivalence}
दो आव्यूह $B, C \in M_{n,k}(A)$ \emph{तुल्य} कहलाते हैं यदि $Q \in GL_n(A)$, $P \in GL_k(A)$
(\emph{$A$ पर} व्युत्क्रमणीय: सारणिक $A^\times$ में) के साथ $C =
QBP$
हो। तुल्य प्रस्तुति आव्यूह तुल्याकारी \omterm{def:b3:modules:module}{मॉड्यूल} $A^n/
\operatorname{im}$ परिभाषित
करते हैं ($A^n$ और $A^k$ में आधार बदलिए)।
\end{definition}

\begin{theorem}[स्मिथ प्रसामान्य रूप]\label{thm:b3:modules:smith}
मान लीजिए $A$ एक \omterm{def:b3:rings:pidufd}{PID} है और $B \in M_{n,k}(A)$। तब $B$ एक विकर्ण आव्यूह के
तुल्य है\index{स्मिथ प्रसामान्य रूप}
\[
\operatorname{diag}(d_1, d_2, \dots, d_r, 0, \dots, 0),
\qquad d_1 \mid d_2 \mid \cdots \mid d_r \neq 0 ,
\]
और $d_i$ सहचारियों तक अद्वितीय हैं: $d_1 \cdots d_i$ $B$ के $i \times i$
उपसारणिकों का महत्तम समापवर्तक है (विशेष रूप से वह समापवर्तक
तुल्यता का एक अपरिवर्ती है)। $d_i$ $B$ के \emph{अपरिवर्ती
गुणनखंड}\index{अपरिवर्ती गुणनखंड} कहलाते हैं।
\end{theorem}

\begin{proof}
\emph{अस्तित्व।} यदि $B = 0$ हो, तो काम पूरा। अन्यथा $B$ के तुल्य
आव्यूहों की प्रविष्टियों से जनित आदर्शों $(b)$ के समुच्चय पर
विचार कीजिए; चूँकि $A$ \omterm{def:b3:rings:noetherian}{नोएथरीय} है, $B$ के तुल्य ऐसा आव्यूह
$B'$ और $B'$ की ऐसी प्रविष्टि $d$ चुनिए जिसके लिए $(d)$ इस
समुच्चय में \emph{महत्तम} हो। पंक्ति और स्तंभ अदला-बदली से $d$
को स्थान $(1,1)$ पर लाइए।

\emph{दावा: $d$ $B'$ की प्रत्येक प्रविष्टि को विभाजित करता
है।} पहले स्तंभ 1: यदि $b_{i1}$ $d$ का गुणज न हो, तो $e = \gcd(d, b_{i1}) = ud +
vb_{i1}$
(बेज़ू) लीजिए, अतः $(e) \supsetneq (d)$। अब $2 \times 2$ आव्यूह की युक्ति: पंक्तियों
$1$ और $i$ पर
\[
\begin{pmatrix} u & v \\ -\dfrac{b_{i1}}e & \dfrac de
\end{pmatrix}
\in GL_2(A)
\qquad
\Bigl(\det = \tfrac{ud + v b_{i1}}e = 1\Bigr)
\]
द्वारा क्रिया करने से एक तुल्य आव्यूह बनता है जिसकी स्थान $(1,1)$
पर प्रविष्टि $e$ है: यह $(d)$ की महत्तमता का खंडन है। अतः
$d$ स्तंभ $1$ को विभाजित करता है, और सममित रूप से पंक्ति
$1$ को भी। पंक्ति 1 और स्तंभ 1 के गुणज घटाने पर वे साफ हो जाते
हैं: $B'$ $\begin{pmatrix} d &
0\\ 0 & B''\end{pmatrix}$ के तुल्य है। इसके बाद, $d$ $B''$ की प्रत्येक
प्रविष्टि $b$ को विभाजित करता है: $b$ की पंक्ति को पंक्ति 1
में जोड़िए (एक प्रारंभिक संक्रिया; नई पहली पंक्ति में $d$ और
$b$-प्रविष्टियाँ हैं) और स्तंभ साफ करने वाला तर्क दोहराइए: कोई
अगुणज फिर से $(d)$ को बेहतर बना देता। अब आकार पर आगमन कीजिए:
$B''$, जिसकी सभी प्रविष्टियाँ $d$ से विभाज्य हैं, का कोई
स्मिथ रूप $\operatorname{diag}(d_2,
\dots)$ है जिसकी प्रविष्टियाँ $d$ से विभाज्य बनी रहती
हैं (किसी भी $QB''P$ की प्रत्येक प्रविष्टि $B''$ की प्रविष्टियों
का $A$-संचय है); अब $d_1 =
d$ रखिए।

\emph{अद्वितीयता।} मान लीजिए $D_i(B)$ समस्त $i \times i$ उपसारणिकों का
महत्तम समापवर्तक है। पंक्ति और स्तंभ संक्रियाएँ, और अधिक व्यापक
रूप से \emph{किसी भी} आव्यूह से गुणन, इस समापवर्तक को छोटा नहीं
कर सकते: $QB$ के $i \times i$ उपसारणिक $B$ के उपसारणिकों के
$A$-संचय होते हैं (कोशी--बिने प्रसार; अथवा सीधे: $QB$ की
प्रत्येक पंक्ति $B$ की पंक्तियों का संचय है, और उपसारणिक
पंक्तियों में बहुरैखिक होते हैं)। अतः $D_i(QBP)$ और $D_i(B)$ परस्पर
विभाजित करते हैं: $D_i$ एक तुल्यता अपरिवर्ती है। विकर्ण रूप पर
शून्येतर $i
\times i$ उपसारणिक $d_j$ में से $i$ के गुणनफल हैं, और
विभाज्यता $d_1 \mid \dots \mid d_r$ से $d_1 \cdots d_i$ ही महत्तम समापवर्तक बन जाता है। अतः
इकाइयों तक $d_1 \cdots d_i = D_i(B)$, और $d_i =
D_i/D_{i-1}$ निर्धारित हो जाता है।
\end{proof}

\begin{method}\label{met:b3:modules:smith}
किसी \omterm{def:b3:rings:pidufd}{यूक्लिडीय प्रांत} ($\Z$, $K[X]$) पर स्मिथ अपचयन एक
कलनविधि है --- यहाँ किसी महत्तमता तर्क की आवश्यकता नहीं: सबसे
छोटे यूक्लिडीय आकार वाली प्रविष्टि को स्थान $(1,1)$ पर लाइए; यदि
वह अपनी पंक्ति या स्तंभ की किसी प्रविष्टि को विभाजित न करे, तो
यूक्लिडीय विभाजन वहाँ कड़ाई से छोटा शेषफल छोड़ देता है --- उसे
भीतर लाकर फिर से आरंभ कीजिए (समाप्ति: आकार घटते जाते हैं); जब वह
अपनी पूरी पंक्ति और स्तंभ को विभाजित करने लगे, तब उन्हें साफ
कीजिए; यदि वह किसी भीतरी प्रविष्टि को विभाजित न करे, तो उस
पंक्ति को पंक्ति $1$ में जोड़कर फिर से आरंभ कीजिए; और भीतरी
खंड पर पुनरावर्तन कीजिए। पूर्णांक आव्यूहों पर व्यवहार में:
प्रविष्टियों का $D_1 = \gcd$, फिर $D_2$, \dots\ छोटे आकारों के लिए
उपसारणिकों से परिकलित कीजिए, अथवा कलनविधि चलाइए।
\end{method}

\begin{example}[एक पूर्ण स्मिथ अपचयन]
\label{ex:b3:modules:smithworked}
$\Z$ पर $M = \begin{pmatrix} 1 & 2 & 3\\ 4 & 5 & 6\\ 7 & 8 & 9
\end{pmatrix}$ का अपचयन कीजिए। कोने की प्रविष्टि $1$ सब कुछ
विभाजित करती है: उसकी पंक्ति और स्तंभ साफ कीजिए ($L_2 \leftarrow L_2 - 4L_1$, $L_3
\leftarrow L_3 - 7L_1$,
फिर $C_2 \leftarrow C_2 - 2C_1$, $C_3
\leftarrow C_3 - 3C_1$):
\[
M \sim \begin{pmatrix} 1 & 0 & 0\\ 0 & -3 & -6\\ 0 & -6 & -12
\end{pmatrix} .
\]
भीतरी खंड में कोने की प्रविष्टि $-3$ सभी प्रविष्टियों को
विभाजित करती है: $L_3
\leftarrow L_3 - 2L_2$ और $C_3 \leftarrow C_3 - 2C_2$ उसे $\operatorname{diag}(-3, 0)$ तक साफ कर देते हैं।
चिह्न समायोजित करने पर (किसी पंक्ति को $-1$ से गुणा कीजिए, जो
वैध संक्रिया है):
\[
M \sim \operatorname{diag}(1, 3, 0),
\qquad
\Z^3/M\Z^3 \cong \Z/3\Z \times \Z .
\]
सारणिक भाजकों से जाँच: $D_1 =
\gcd(\text{प्रविष्टियाँ}) = 1$; $M$ का प्रत्येक $2\times2$ उपसारणिक
$3$ का गुणज है (उदाहरणार्थ $\det\bigl(\begin{smallmatrix}1 & 2\\ 4
& 5\end{smallmatrix}\bigr) = -3$) और एक $-3$ के बराबर है:
$D_2 =
3$; $D_3 = \det M = 0$। अतः $d_1 = 1$, $d_2 = 3$, $d_3 = 0$: वही उत्तर। दो शिक्षाएँ:
शून्य \omterm{thm:b3:modules:smith}{अपरिवर्ती गुणनखंड} कोटि की गिरावट दर्ज करता है (सह-अष्टि एक
\omterm{def:b3:modules:free}{मुक्त} $\Z$ योज्यांश उठा लेती है), और विभाज्यता शृंखला $1 \mid 3 \mid 0$ ही
स्मिथ का प्रमाणपत्र है --- ऐसा विकर्ण अपचयन जो इस शृंखला का
उल्लंघन करता हो (मान लीजिए $\operatorname{diag}(2, 3)$, जिसे असावधान पाठक $\bigl(\begin{smallmatrix}2 & 0\\ 0 &
3\end{smallmatrix}\bigr)$ से
बहुत जल्दी रुककर उत्पन्न कर सकता है: सही स्मिथ रूप $\operatorname{diag}(1, 6)$ है,
क्योंकि यहाँ $D_1 = 1$!) अभी पूरा नहीं हुआ है।
\end{example}

\begin{omfigure}
\begin{tikzpicture}[scale=0.72]
  \clip (-3.4,-1.3) rectangle (4.4,4.3);
  \foreach \x in {-3,...,4}
    \foreach \y in {-1,...,4}
      \fill[black!40] (\x,\y) circle (1.6pt);
  \fill[omDef!12, opacity=0.7] (0,0) -- (1,3) -- (3,3) -- (2,0)
    -- cycle;
  \foreach \a/\b in {0/0, 2/0, 4/0, -2/0, 1/3, 3/3, -1/3, -3/3,
                     2/6, -2/-6, 0/6}
    \fill[omThm] (\a,\b) circle (2.6pt);
  \draw[omThm, thick, ->] (0,0) -- (2,0)
    node[midway, below, font=\small] {$(2,0)$};
  \draw[omThm, thick, ->] (0,0) -- (1,3)
    node[midway, left, font=\small] {$(1,3)$};
\end{tikzpicture}

{\small $\Z^2$ का उपजालक $L = \Z(2,0) + \Z(1,3)$ (लाल बिंदु)। $\bigl(\begin{smallmatrix} 2 &
1\\ 0 & 3\end{smallmatrix}\bigr)$ का \omterm{thm:b3:modules:smith}{स्मिथ
प्रसामान्य रूप} $\operatorname{diag}(1, 6)$ है: $\Z^2$ के अनुकूलित आधार $f_1 = (1,3)$, $f_2 = (0,1)$ में
$L = \Z f_1 \oplus \Z\,6f_2$ होता है, अतः $\Z^2/L \cong \Z/6\Z$ --- सूचकांक $\abs{\det} = 6$ के बराबर है, अर्थात्
छायांकित मूल प्रांत के क्षेत्रफल के।}
\end{omfigure}

\section{संरचना प्रमेय}

\begin{definition}\label{def:b3:modules:torsion}
मान लीजिए $A$ एक प्रांत है और $M$ एक $A$-मॉड्यूल।
\emph{मरोड़}\index{मरोड़} उपमॉड्यूल है
\[
T(M) = \{x \in M : ax = 0 \text{ किसी } a \neq 0 \text{ के लिए}\}
\]
(यह उपमॉड्यूल है: यदि $ax = by = 0$ हो तो $ab(x + y) = 0$, $ab \ne 0$)। $M$
\emph{मरोड़-मुक्त} कहलाता है यदि $T(M) = 0$ हो, और \emph{मरोड़
\omterm{def:b3:modules:module}{मॉड्यूल}} यदि $T(M) = M$ हो।
\end{definition}

\begin{theorem}[किसी PID पर परिमिततः जनित मॉड्यूलों की
संरचना]\label{thm:b3:modules:structure}
मान लीजिए $A$ एक \omterm{def:b3:rings:pidufd}{PID} है और $M$ एक \omterm{def:b3:modules:free}{परिमिततः जनित}
$A$-मॉड्यूल। तब एक अद्वितीय $r \in \N$ और सहचारियों तक अद्वितीय
शून्येतर अनैकिक अवयव $d_1 \mid d_2 \mid
\cdots \mid d_s$ ऐसे विद्यमान हैं
कि\index{मॉड्यूलों के लिए संरचना प्रमेय}
\[
M \;\cong\; A^r \,\oplus\, A/(d_1) \oplus \cdots \oplus A/(d_s).
\]
इसके अतिरिक्त $T(M) \cong A/(d_1)\oplus\dots\oplus A/(d_s)$ और $M/T(M)
\cong A^r$: किसी \omterm{def:b3:rings:pidufd}{PID} पर \omterm{def:b3:modules:free}{परिमिततः जनित}
\emph{\omterm{def:b3:modules:torsion}{मरोड़-मुक्त}} \omterm{def:b3:modules:module}{मॉड्यूल} \omterm{def:b3:modules:free}{मुक्त} होता है।
\end{theorem}

\begin{proof}
\emph{अस्तित्व।} $M \cong A^n/\operatorname{im}(C)$ प्रस्तुत कीजिए (\cref{rem:b3:modules:presentation}) और $C$ को स्मिथ
रूप में लाइए: दोनों आधार-परिवर्तनों के बाद $M \cong A^n / (d_1A \times \dots
\times d_rA \times 0 \times \dots \times 0) \cong A/(d_1) \oplus
\dots \oplus A/(d_r) \oplus A^{\,n - r}$। जिन गुणनखंडों
में $d_i$ इकाई है उन्हें छोड़ दीजिए ($A/(d_i) = 0$); विभाज्यता शृंखला
बनी रहती है।

\emph{\omterm{def:b3:modules:torsion}{मरोड़} की पहचान।} इस विघटन में $A^r$ \omterm{def:b3:modules:torsion}{मरोड़-मुक्त} है (प्रांत
में कोई शून्य भाजक नहीं) और प्रत्येक $A/(d_i)$ \omterm{def:b3:modules:torsion}{मरोड़} है ($d_i \ne 0$ से
नष्ट); प्रत्यक्ष योग \omterm{def:b3:modules:torsion}{मरोड़} को उसी के अनुसार विभाजित कर देता है:
$T(M) = \bigoplus_i A/(d_i)$ और $M/T(M) \cong A^r$।

\emph{$r$ की अद्वितीयता}: $M/T(M) \cong A^r$ केवल $M$ पर निर्भर करता है,
और \cref{prop:b3:modules:rank} से $r$ निश्चित हो जाता है।

\emph{$d_i$ की अद्वितीयता}: केवल \omterm{def:b3:modules:torsion}{मरोड़ मॉड्यूल} $T = T(M)$ की स्थिति
देखनी पर्याप्त है। प्रत्येक $d_i$ को अभाज्यों में विघटित कीजिए
और चीनी शेषफल प्रमेय से विभाजित कीजिए (\cref{thm:b3:rings:crt}; भिन्न अभाज्य
\omterm{thm:b3:rings:crt}{सह-महत्तम आदर्श} जनित करते हैं):
\[
A/(d) \cong \bigoplus_{p \mid d} A/\bigl(p^{v_p(d)}\bigr):
\qquad
T \cong \bigoplus_{p} \bigoplus_{j} A/\bigl(p^{k_{p,j}}\bigr),
\]
और \emph{प्रारंभिक भाजक}\index{प्रारंभिक भाजक} $p^{k_{p,j}}$ मिलते हैं।
विलोमतः $d_i$ को \omterm{thm:b3:modules:structure}{प्रारंभिक भाजकों} के बहुसमुच्चय से पुनर्निर्मित
किया जा सकता है ($d_s$ = प्रत्येक अभाज्य की उच्चतम घात का
गुणनफल, इत्यादि), अतः केवल यह सिद्ध करना पर्याप्त है कि प्रत्येक
अभाज्य $p$ के लिए बहुसमुच्चय $\{k_{p,j}\}_j$ $T$ से निर्धारित हो जाता
है। $p$ स्थिर कीजिए; $j \geq 1$ के लिए $A/(p)$-सदिश समष्टियों $p^{j-1}T/p^jT$
पर विचार कीजिए। किसी चक्रीय गुणनखंड $A/(p^k)$ पर:
\[
p^{j-1}\bigl(A/(p^k)\bigr)\big/p^{j}\bigl(A/(p^k)\bigr)
\cong
\begin{cases}
A/(p) & \text{यदि } j \leq k,\\
0 & \text{यदि } j > k,
\end{cases}
\]
और किसी गुणनखंड $A/(q^k)$, $q \neq p$, पर $p$ से गुणन वहाँ एकैकी आच्छादक
है ($q^k$ के सापेक्ष $p$ व्युत्क्रमणीय है: बेज़ू), अतः भागफल
$0$ है। प्रत्यक्ष योग इनसे होकर गुजरते हैं: $\dim_{A/(p)}
p^{j-1}T/p^jT = \#\{i : k_{p,i} \geq j\}$। ये
अंतर्जात विमाएँ घातांकों के बहुसमुच्चय को निर्धारित कर देती हैं।
\end{proof}

\begin{corollary}[परिमिततः जनित आबेली समूह]
\label{cor:b3:modules:abelian}
प्रत्येक \omterm{def:b3:modules:free}{परिमिततः जनित} आबेली समूह $d_1 \mid \dots \mid
d_s$ के साथ अद्वितीय रूप से
$\Z^r \times
\Z/d_1\Z\times\dots\times\Z/d_s\Z$ होता है। प्रत्येक \emph{परिमित} आबेली समूह अभाज्य-घात क्रम
के चक्रीय समूहों का गुणन है, जो बहुसमुच्चय के रूप में अद्वितीय
है।\index{आबेली समूह, संरचना}
\end{corollary}

\begin{example}\label{ex:b3:modules:abelian}
क्रम $p^n$ के आबेली समूह $n$ के \emph{विभाजनों} के अनुरूप हैं:
$p^4$ के लिए $\Z/p^4$, $\Z/p^3\times\Z/p$, $(\Z/p^2)^2$, $\Z/p^2 \times (\Z/p)^2$, $(\Z/p)^4$ --- पाँच समूह,
क्योंकि $4$ के पाँच विभाजन हैं। मिश्रित क्रमों में गणनाएँ
अभाज्य-दर-अभाज्य गुणा हो जाती हैं (चीनी शेषफल प्रमेय): क्रम
$2^4 \cdot 3^2 = 144$ के $5 \times 2$ आबेली समूह हैं।
\end{example}

\section{अनुप्रयोग: अंतःसमाकारिताओं के विहित रूप}

मान लीजिए $K$ एक क्षेत्र है, $V$ परिमित विमा $n$ की एक
$K$-सदिश समष्टि, और $u \in \mathcal L(V)$; $P
\cdot x = P(u)(x)$ (\cref{ex:b3:modules:examples}) द्वारा $V$ को
$K[X]$-मॉड्यूल बना दीजिए। यह \omterm{def:b3:modules:module}{मॉड्यूल} \omterm{def:b3:modules:free}{परिमिततः जनित} है (कोई भी
$K$-आधार उसे जनित करता है) और \omterm{def:b3:modules:torsion}{मरोड़} है: प्रत्येक $x$ के लिए
$n+1$ सदिश $x, u(x), \dots, u^n(x)$ $K$-आश्रित होते हैं, जिससे एक शून्येतर
विनाशक बहुपद मिल जाता है।

\begin{definition}\label{def:b3:modules:companion}
मोनिक $P = X^m + a_{m-1}X^{m-1} + \dots + a_0$ के लिए \emph{सहचर आव्यूह}\index{सहचर आव्यूह} है
\[
C_P = \begin{pmatrix}
0 & & & -a_0\\
1 & \ddots & & -a_1\\
& \ddots & 0 & \vdots\\
& & 1 & -a_{m-1}
\end{pmatrix} :
\]
अर्थात् आधार $1, \bar X, \dots, \bar X^{m-1}$ में $K[X]/(P)$ पर ``$X$ से गुणन'' का आव्यूह।
\end{definition}

\begin{theorem}[फ्रोबेनियस: परिमेय विहित रूप]
\label{thm:b3:modules:frobenius}
मोनिक अनचर बहुपदों का एक अद्वितीय अनुक्रम $P_1
\mid P_2 \mid \dots \mid P_s$ ($u$ के
\emph{समरूपता अपरिवर्ती}\index{समरूपता अपरिवर्ती}) ऐसा विद्यमान
है कि $K[X]$-मॉड्यूलों के रूप में
\[
V \cong K[X]/(P_1) \oplus \cdots \oplus K[X]/(P_s):
\]
और किसी उपयुक्त आधार में $u$ का आव्यूह खंड-विकर्ण $\operatorname{diag}(C_{P_1}, \dots, C_{P_s})$ होता
है। इसके अतिरिक्त:
\begin{enumerate}
  \item $P_s = \mu_u$ (न्यूनतम बहुपद) और $P_1\cdots P_s =
        \chi_u$ (अभिलक्षणिक बहुपद);
        विशेष रूप से $\mu_u
        \mid \chi_u$ \emph{(केली--हैमिल्टन की पुनः
        उपपत्ति)} तथा $\chi_u \mid \mu_u^{\,s}$, अतः $\chi_u$ और $\mu_u$ के \omterm{def:b3:rings:divisibility}{अखंडनीय}
        गुणनखंड एक ही हैं।
  \item दो अंतःसमाकारिताएँ (या वर्ग आव्यूह) समरूप हैं तभी और
        केवल तभी जब उनके \omterm{thm:b3:modules:frobenius}{समरूपता अपरिवर्ती} एक ही हों।
\end{enumerate}
\end{theorem}

\begin{proof}
\omterm{def:b3:rings:pidufd}{PID} $K[X]$ पर संरचना प्रमेय (\cref{thm:b3:modules:structure}) लगाइए: \omterm{def:b3:modules:torsion}{मरोड़ मॉड्यूल} $V$
अपरिवर्ती गुणनखंडों $P_i$ के साथ विघटित होता है, जिन्हें मोनिक
मानकीकृत किया गया है ($K[X]$ की इकाइयाँ $K^\times$ हैं); कोई \omterm{def:b3:modules:free}{मुक्त}
भाग नहीं आता ($V$ \omterm{def:b3:modules:torsion}{मरोड़} है)। प्रत्येक चक्रीय गुणनखंड $K[X]/(P_i)$ पर
$X$ से गुणन का आव्यूह $\bar X$ की घातों के आधार में $C_{P_i}$ है:
आधारों को जोड़ने पर खंड रूप मिल जाता है।

(1) $V = \bigoplus K[X]/(P_i)$ का विनाशक $(P_1)\cap
\dots\cap(P_s) = (P_s)$ है (विभाज्यता शृंखला: $P_s$ एक
उभयनिष्ठ गुणज है, और अंतिम गुणनखंड में $1$ का वर्ग ठीक $(P_s)$
से ही नष्ट होता है): अतः $\mu_u = P_s$। $\chi_u$ के लिए: किसी चक्रीय गुणनखंड
पर $m = \deg P$ पर आगमन द्वारा $\chi_{C_P} = P$। $\det(XI_m - C_P)$ को पहली पंक्ति के अनुदिश
प्रसारित कीजिए (जिसकी प्रविष्टियाँ $X$, फिर शून्य, फिर अंतिम
स्तंभ में $a_0$ हैं):
\[
\det(XI_m - C_P) = X\,\det\bigl(XI_{m-1} - C_{\tilde P}\bigr)
+ (-1)^{1+m}\,a_0\,\det L ,
\]
जहाँ $\tilde P = X^{m-1} + a_{m-1}X^{m-2} + \dots + a_1$ (वही आकृति, एक आकार छोटी) और $L$ त्रिभुजाकार है
जिसका विकर्ण $(-1, \dots, -1)$ है, अतः $\det L = (-1)^{m-1}$। आगमन से पहला पद $X\tilde P$ है और
दूसरा $a_0$: कुल मिलाकर $X\tilde P + a_0 = P$ (आधार स्थिति $m=1$: $\det(X + a_0) = P$)।
सारणिक खंडों पर गुणा होते हैं: $\chi_u = \prod P_i$। केली--हैमिल्टन: $\chi_u \in (\mu_u)$,
क्योंकि प्रत्येक $P_s \mid$ \dots विलोमतः प्रत्येक $P_i \mid P_s$, अतः $\chi_u = \prod P_i$
$P_s^{\,s} = \mu_u^s$ को विभाजित करता है; और $\mu_u =
P_s$ $\chi_u$ को उसके एक गुणनखंड के
रूप में विभाजित करता है।

(2) समरूप अंतःसमाकारिताएँ संयुग्म \omterm{def:b3:modules:module}{मॉड्यूल} संरचनाएँ हैं, अतः उनके
अपरिवर्ती समान होते हैं (\cref{thm:b3:modules:structure} में अद्वितीयता); विलोमतः समान
अपरिवर्ती तुल्याकारी $K[X]$-मॉड्यूल देते हैं, और \omterm{def:b3:modules:module}{मॉड्यूल}
तुल्याकारिता ठीक वही रैखिक एकैकी आच्छादक प्रतिचित्रण है जो दोनों
अंतःसमाकारिताओं को आपस में गूँथता है: अर्थात् एक समरूपता।
\end{proof}

\begin{corollary}[समरूपता क्षेत्र विस्तार के प्रति असंवेदी है]
\label{cor:b3:modules:descent}
मान लीजिए $K \subseteq L$ क्षेत्र हैं और $M, N \in M_n(K)$। यदि $M$ और $N$ $L$
पर समरूप हों, तो वे $K$ पर भी समरूप हैं।
\end{corollary}

\begin{proof}
$M$ के \omterm{thm:b3:modules:frobenius}{समरूपता अपरिवर्ती} $K[X]$ पर प्रस्तुति आव्यूह $XI_n - M$ पर
स्मिथ के उपसारणिक सूत्र (\cref{thm:b3:modules:smith}) से परिकलित होते हैं --- वस्तुतः
$K[X]$-मॉड्यूल $V_M = K^n$ की प्रस्तुति $XI_n - M$ है: प्रतिचित्रण $K[X]^n \to V_M$,
$(Q_i) \mapsto \sum Q_i(M)e_i$, आच्छादक है और उसकी अष्टि $XI_n - M$ के स्तंभों से जनित है
(सीधा सत्यापन: उन स्तंभों के सापेक्ष $K[X]^n$ का प्रत्येक अवयव किसी
अचर सदिश तक अपचयित हो जाता है, और अचर सदिश एकैकी आच्छादक रूप से
जाते हैं; सप्ताहांत समस्या इसे विस्तार से लिखती है)। बहुपदों के
महत्तम समापवर्तक क्षेत्र विस्तार से नहीं बदलते: यदि $d$ किसी
कुल $(f_j)$ का $K[X]$ में मोनिक महत्तम समापवर्तक हो, तो बेज़ू से
$u_j \in K[X]$ के साथ $d = \sum u_jf_j$ मिलता है, अतः \emph{$L[X]$ में} $f_j$ का
प्रत्येक उभयनिष्ठ भाजक $d$ को विभाजित करता है; और चूँकि $d$
स्वयं एक उभयनिष्ठ भाजक है, वह $L[X]$ में भी महत्तम समापवर्तक है।
अतः $XI_n - M$ के \omterm{thm:b3:modules:smith}{अपरिवर्ती गुणनखंड}, जो क्रमागत उपसारणिक समापवर्तकों
के भागफल हैं, $K$ पर और $L$ पर एक ही रहते हैं: अर्थात्
$M, N$ के $L$ पर \omterm{thm:b3:modules:frobenius}{समरूपता अपरिवर्ती} वही हैं जो $K$ पर; अब
\cref{thm:b3:modules:frobenius}(2) से निष्कर्ष निकालिए।
\end{proof}

\begin{theorem}[जॉर्डन रूप, पुनः व्युत्पन्न]\label{thm:b3:modules:jordan}
मान लीजिए $\chi_u$ $K$ पर विभाजित हो जाता है (उदाहरणार्थ $K = \C$)।
अपरिवर्ती गुणनखंडों के स्थान पर $V$ पर प्रारंभिक-भाजक विघटन
(\cref{thm:b3:modules:structure} की उपपत्ति) लगाने पर
\[
V \cong \bigoplus_{\lambda, j} K[X]\big/\bigl((X -
\lambda)^{k_{\lambda,j}}\bigr),
\]
और प्रत्येक गुणनखंड के आधार $\bigl(\overline{(X-\lambda)^{k-1}}, \dots,
\overline{(X - \lambda)}, \bar 1\bigr)$ में $u$ जॉर्डन खंड $J_k(\lambda)$ की
तरह क्रिया करता है: अर्थात् विभाजित अभिलक्षणिक बहुपद वाली
प्रत्येक अंतःसमाकारिता का कोई जॉर्डन आधार होता है, और खंडों का
बहुसमुच्चय $(\lambda, k)$ अद्वितीय है।\index{जॉर्डन रूप}
\end{theorem}

\begin{proof}
\omterm{def:b3:modules:torsion}{मरोड़ मॉड्यूल} $V$ के \omterm{thm:b3:modules:structure}{प्रारंभिक भाजक} $(X - \lambda)^k$ हैं, जहाँ $X - \lambda$
$\mu_u$ के \omterm{def:b3:rings:divisibility}{अखंडनीय} गुणनखंडों में विचरता है (और वह विभाजित हो
जाता है, क्योंकि $\chi_u$ विभाजित होता है और दोनों के \omterm{def:b3:rings:divisibility}{अखंडनीय}
गुणनखंड एक ही हैं, \cref{thm:b3:modules:frobenius})। $W = K[X]/((X-\lambda)^k)$ में $j = 1, \dots, k$ के लिए $f_j = \overline{(X - \lambda)^{k-j}}$ रखिए:
तब $(X - \lambda)f_j = f_{j-1}$ (जहाँ $f_0 = 0$), अर्थात् $u(f_j)
= \lambda f_j + f_{j-1}$: $(f_1, \dots,
f_k)$ पर $u$ का आव्यूह
ठीक $J_k(\lambda)$ है (विकर्ण के ऊपर एक)। \omterm{thm:b3:modules:structure}{प्रारंभिक भाजकों} के बहुसमुच्चय
की अद्वितीयता \cref{thm:b3:modules:structure} है।
\end{proof}

\begin{remark}\label{rem:b3:modules:overview}
विहित रूपों का क्रम अब पारदर्शी है: \emph{परिमेय} रूप प्रत्येक
क्षेत्र पर विद्यमान है और समरूपता को निरपेक्ष रूप से पकड़ लेता
है (\cref{cor:b3:modules:descent}); \emph{जॉर्डन} रूप उसका परिष्कार है, जब $\chi_u$ विभाजित
हो जाता है। दूसरे वर्ष की विमा-गणना वाली जॉर्डन उपपत्तियाँ इसी
में समा जाती हैं: वह सारी संचयिकी वस्तुतः \omterm{def:b3:rings:pidufd}{PID} $K[X]$ का अंकगणित
थी।
\end{remark}

\section{अभ्यास}

\begin{exercise}[$\star$]\label{exo:b3:modules:1}
(क) दर्शाइए कि $\Q$ $\Z$-मॉड्यूल के रूप में \omterm{def:b3:modules:free}{परिमिततः जनित} नहीं
है।
(ख) दर्शाइए कि $\Q$ \omterm{def:b3:modules:torsion}{मरोड़-मुक्त} है पर \omterm{def:b3:modules:free}{मुक्त} नहीं।
(ग) दोनों में से कोई भी कथन \cref{thm:b3:modules:structure} का खंडन क्यों नहीं करता?
\end{exercise}

\begin{exercise}[$\star$]\label{exo:b3:modules:2}
तुल्याकारिता तक क्रम $360$ के आबेली समूहों को प्रारंभिक-भाजक और
अपरिवर्ती-गुणनखंड दोनों रूपों में सूचीबद्ध कीजिए। क्रम $p^5$ के
कितने आबेली समूह हैं?
\end{exercise}

\begin{exercise}[$\star$]\label{exo:b3:modules:3}
$\Z$ पर
\[
B = \begin{pmatrix} 2 & 4\\ 6 & 8 \end{pmatrix},
\qquad
C = \begin{pmatrix} 2 & 0 & 0\\ 0 & 3 & 0\\ 0 & 0 & 12
\end{pmatrix},
\]
का \omterm{thm:b3:modules:smith}{स्मिथ प्रसामान्य रूप} परिकलित कीजिए, और आबेली समूहों $\Z^2/B\Z^2$
तथा $\Z^3/C\Z^3$ की पहचान कीजिए।
\end{exercise}

\begin{exercise}[$\star\star$]\label{exo:b3:modules:4}
मान लीजिए $L \subseteq \Z^n$ कोटि $n$ का उपसमूह है, जिसका आधार $B \in M_n(\Z)$,
$\det B \neq 0$, के स्तंभ हैं। दर्शाइए कि $\Z^n/L$ परिमित है और उसकी गणनीयता
$\abs{\det B}$ है, तथा $B$ के अपरिवर्ती गुणनखंडों $d_i$ के लिए $\Z^n/L
\cong \prod_i \Z/d_i\Z$।
$L = \Z(2,0) + \Z(1,3)$ से उदाहरण दीजिए।
\end{exercise}

\begin{exercise}[$\star\star$]\label{exo:b3:modules:5}
मान लीजिए $A$ एक प्रांत है। (क) सत्यापित कीजिए कि $T(M)$ एक
उपमॉड्यूल है और $M/T(M)$ \omterm{def:b3:modules:torsion}{मरोड़-मुक्त} है। (ख) दर्शाइए कि $K[X,Y]$ का
\omterm{def:b3:rings:ideal}{आदर्श} $(X, Y)$, $K[X,Y]$-मॉड्यूल के रूप में, \omterm{def:b3:modules:torsion}{मरोड़-मुक्त} है पर \omterm{def:b3:modules:free}{मुक्त}
नहीं: संरचना प्रमेय को \omterm{def:b3:rings:pidufd}{PID} परिकल्पना की वास्तव में आवश्यकता है।
\end{exercise}

\begin{exercise}[$\star\star$]\label{exo:b3:modules:6}
(क) दर्शाइए कि $\Z$-मॉड्यूल $\Z$ में $2\Z$ का कोई प्रत्यक्ष
पूरक नहीं है: \omterm{def:b3:modules:free}{मुक्त} मॉड्यूलों के उपमॉड्यूल \omterm{def:b3:modules:free}{मुक्त} होते हैं
(\cref{thm:b3:modules:submodule}), पर उनका प्रत्यक्ष योज्यांश होना आवश्यक नहीं। (ख)
दर्शाइए कि यदि $M \subseteq A^n$ ($A$ एक \omterm{def:b3:rings:pidufd}{PID}) इस शर्त को संतुष्ट करे कि
$A^n/M$ \omterm{def:b3:modules:torsion}{मरोड़-मुक्त} है, तो $M$ प्रत्यक्ष योज्यांश \emph{होता}
है।
\end{exercise}

\begin{exercise}[$\star\star$]\label{exo:b3:modules:7}
\cref{exo:b3:modules:3} के स्मिथ रूप का उपयोग करते हुए (दोनों ओर चरों के
व्युत्क्रमणीय परिवर्तन) $\Z^2$ में निकाय
\[
\begin{cases}
2x + 4y \equiv b_1 \pmod{20},\\
6x + 8y \equiv b_2 \pmod{20},
\end{cases}
\]
हल कीजिए: किन युग्मों $(b_1, b_2)$ के लिए हल विद्यमान हैं?
\end{exercise}

\begin{exercise}[$\star\star$]\label{exo:b3:modules:8}
(क) किसी शून्यभावी $4 \times 4$ आव्यूह के समस्त \omterm{thm:b3:modules:frobenius}{समरूपता अपरिवर्ती} और
संभव \omterm{thm:b3:modules:jordan}{जॉर्डन रूप} ज्ञात कीजिए, और उन्हें $4$ के उस विभाजन के
अनुसार क्रमबद्ध कीजिए जिसे वे साकार करते हैं।
(ख) समान अभिलक्षणिक \emph{और} न्यूनतम बहुपद वाले दो ऐसे $4\times4$
सम्मिश्र आव्यूह प्रस्तुत कीजिए जो समरूप न हों, और सिद्ध कीजिए कि
$n \leq 3$ के लिए ऐसा नहीं हो सकता।
\end{exercise}

\begin{exercise}[$\star\star\star$]\label{exo:b3:modules:9}
मान लीजिए $u \in \mathcal L(V)$, $\dim V = n$। दर्शाइए कि निम्नलिखित समतुल्य हैं:
(क) $V$ एक चक्रीय $K[X]$-मॉड्यूल है (ऐसा $x$ है जिसके लिए
$V = K[u]x$ हो --- एक \emph{चक्रीय सदिश}); (ख) $\mu_u = \chi_u$; (ग) \cref{thm:b3:modules:frobenius} में
$s = 1$। इससे निष्कर्ष निकालिए कि किसी \omterm{def:b3:modules:companion}{सहचर आव्यूह} का चक्रीय सदिश
होता है, और निर्धारित कीजिए कि किसी विकर्ण आव्यूह का कब होता है।
\end{exercise}

\begin{exercise}[$\star\star\star$]\label{exo:b3:modules:10}
$\Z^n$ की अंतःसमाकारिता के रूप में देखे गए $M \in M_n(\Z)$ के लिए सूचकांक
सूत्र सिद्ध कीजिए: यदि $\det M \ne 0$ हो, तो $[\Z^n : M\Z^n] = \abs{\det M}$, और इससे निष्कर्ष
निकालिए कि $M \in GL_n(\Z)$ तभी और केवल तभी जब $\det M = \pm 1$। अनुप्रयोग: समूह
$\Z^2$ के सूचकांक $m$ के ठीक $\sigma_1(m) = \sum_{d \mid m} d$ उपसमूह हैं। \emph{(एर्मीत
रूप $\bigl(\begin{smallmatrix}
a & b\\ 0 & d\end{smallmatrix}\bigr)$, $ad = m$, $0 \leq b <
d$ के आव्यूह गिनिए।)}
\end{exercise}

\begin{exercise}[$\star\star$]\label{exo:b3:modules:11}
(आबेली समूहों में समीकरण $x^k = e$) मान लीजिए $G$ एक परिमित आबेली
समूह है जिसके \omterm{thm:b3:modules:smith}{अपरिवर्ती गुणनखंड} $d_1 \mid d_2 \mid \dots
\mid d_s$ हैं।
(क) दर्शाइए कि प्रत्येक $k \geq 1$ के लिए
\[
\#\{x \in G : x^k = e\} \;=\; \prod_{i=1}^{s}\gcd(k, d_i) .
\]
(ख) निष्कर्ष निकालिए: कोई परिमित आबेली समूह चक्रीय है तभी और
केवल तभी जब प्रत्येक $k$ के लिए समीकरण $x^k = e$ के अधिकतम $k$
हल हों।
(ग) $K^\times$ ($K$ एक क्षेत्र, \cref{ch:b3:galois}) के परिमित उपसमूहों की
चक्रीयता पुनः प्राप्त कीजिए: बहुपद $X^k - 1$ (ख) की कसौटी की गारंटी
क्यों देता है?
\end{exercise}

\begin{exercise}[$\star\star\star$]\label{exo:b3:modules:12}
(प्रारंभिक आव्यूह जनक हैं) (क) दर्शाइए कि $M \in M_n(\Z)$ $M_n(\Z)$ में
व्युत्क्रमणीय है तभी और केवल तभी जब $\det M = \pm1$।
(ख) दर्शाइए कि $SL_2(\Z)$ दो प्रारंभिक आव्यूहों
$E = \bigl(\begin{smallmatrix}1 & 1\\ 0 &
1\end{smallmatrix}\bigr)$ और
$F = \bigl(\begin{smallmatrix}1 & 0\\ 1 &
1\end{smallmatrix}\bigr)$
से जनित है।
\emph{($E, F$ की घातों से बाएँ गुणन द्वारा $M
\in SL_2(\Z)$ के पहले स्तंभ पर
यूक्लिडीय कलनविधि चलाइए और $\pm\bigl(\begin{smallmatrix}1 & *\\ 0 &
1\end{smallmatrix}\bigr)$ तक पहुँचिए; शेष हाथ से पूरा
कीजिए --- ध्यान दीजिए कि $-I = (EF^{-1}E)^2$।)}
(ग) स्मिथ अपचयन से संबंध समझाइए: $\Z$ पर सारणिक $1$ वाली
पंक्ति और स्तंभ संक्रियाएँ चिह्नों तक विकर्णन के लिए पर्याप्त
हैं।
\end{exercise}

\section{समस्या: क्रमविनिमेयी और द्विक क्रमविनिमेयी}

\begin{problem}[{सप्ताहांत समस्या --- परिमेय रूप, क्रमविनिमेयी,
द्विक क्रमविनिमेयी}]\label{pb:b3:modules:1}
मान लीजिए $K$ एक क्षेत्र है, $V$ विमा $n \geq
1$ की एक
$K$-सदिश समष्टि, और $u \in \mathcal L(V)$। हम \emph{क्रमविनिमेयी}
\[
\mathcal C(u) = \{v \in \mathcal L(V) : uv = vu\},
\]
का अध्ययन करते हैं, जो $K[u] = \{P(u) : P \in
K[X]\}$ को अंतर्विष्ट करने वाला $\mathcal L(V)$ का
उपबीजगणित है, और फ्रोबेनियस का विमा सूत्र तथा द्विक क्रमविनिमेयी
प्रमेय सिद्ध करते हैं: $\mathcal C(\mathcal C(u)) = K[u]$। सर्वत्र $V$ $u$ द्वारा
परिभाषित $K[X]$-मॉड्यूल है, जिसके \omterm{thm:b3:modules:smith}{अपरिवर्ती गुणनखंड} $P_1 \mid \cdots \mid P_s$ हैं और
जिसका चक्रीय विघटन $V =
\bigoplus_{i=1}^s V_i$, $V_i = K[u]\,x_i \cong K[X]/(P_i)$, $n_i
= \deg P_i$ (\cref{thm:b3:modules:frobenius}) है।

\medskip
\noindent\textbf{भाग I --- प्रस्तुति आव्यूह $XI - M$, और प्रारंभिक
अभ्यास।}
\begin{enumerate}
  \item मान लीजिए $M \in M_n(K)$ है और $\varphi \colon K[X]^n \to V_M
        = K^n$ $(Q_1, \dots, Q_n)$ को $\sum_i Q_i(M)e_i$ पर भेजता है।
        दर्शाइए कि $\varphi$ $K[X]$-मॉड्यूलों की आच्छादक समाकारिता है
        और $XI_n - M$ का प्रत्येक स्तंभ $\ker\varphi$ में है।
  \item दर्शाइए कि $XI_n - M$ के स्तंभों के सापेक्ष $K[X]^n$ का प्रत्येक
        अवयव किसी अचर सदिश के सर्वांगसम है \emph{($Xe_i \equiv \sum_j m_{ji}e_j$ का उपयोग
        करके घातें घटाइए)}, और इससे $\ker\varphi = (XI_n - M)\,K[X]^n$ निष्कर्ष निकालिए:
        \omterm{def:b3:modules:module}{मॉड्यूल} $V_M$ का प्रस्तुति आव्यूह $XI_n - M$ है। \cref{cor:b3:modules:descent} का
        प्रस्थान बिंदु पुनः प्राप्त कीजिए: $M$ के \omterm{thm:b3:modules:frobenius}{समरूपता
        अपरिवर्ती} $XI_n - M$ के अनैकिक \omterm{thm:b3:modules:smith}{अपरिवर्ती गुणनखंड} हैं।
  \item इनके \omterm{thm:b3:modules:frobenius}{समरूपता अपरिवर्ती} परिकलित कीजिए: अदिश आव्यूह $\lambda I_n$;
        भिन्न-भिन्न विकर्ण प्रविष्टियों वाला विकर्ण आव्यूह;
        $n \times n$ जॉर्डन खंड $J_n(0)$; और $n = 3$ के लिए $\operatorname{diag}(J_2(0), J_1(0))$।
  \item दर्शाइए कि $\dim K[u] = \deg \mu_u = n_s$।
\end{enumerate}

\medskip
\noindent\textbf{भाग II --- चक्रीय मॉड्यूलों के बीच समाकारिताएँ।}
\begin{enumerate}[resume]
  \item मान लीजिए $P, Q$ मोनिक अनचर हैं। दर्शाइए कि कोई
        $K[X]$-समाकारिता $f \colon K[X]/(P) \to K[X]/(Q)$ $f(\bar 1)$ से निर्धारित हो जाती है, और
        यह कि $\bar c \in K[X]/(Q)$ $f(\bar 1)$ का काम कर सकता है तभी और केवल तभी जब
        $K[X]/(Q)$ में $P\bar c = 0$।
  \item इससे निष्कर्ष निकालिए
        \[
        \operatorname{Hom}_{K[X]}\bigl(K[X]/(P),\, K[X]/(Q)\bigr)
        \;\cong\; K[X]\big/\bigl(\gcd(P, Q)\bigr),
        \]
        जिसकी $K$ पर विमा $\deg \gcd(P, Q)$ है। \emph{(दर्शाइए कि $K[X]/(Q)$ में
        $P\bar c = 0$ के हल $\bar c$ $\overline{Q/\gcd(P,Q)}$ से जनित चक्रीय उपमॉड्यूल बनाते
        हैं।)}
  \item \emph{फ्रोबेनियस का सूत्र} सिद्ध कीजिए:
        \[
        \dim_K \mathcal C(u) = \sum_{i,j=1}^{s} \deg\gcd(P_i,
        P_j) = \sum_{i=1}^{s} (2s - 2i + 1)\, n_i .
        \]
        \emph{(क्रमविनिमेय करने वाला $v$ ठीक $V$ की
        $K[X]$-अंतःसमाकारिता है; $\operatorname{End}(\bigoplus_i V_i)$ को समाकारिताओं $V_j \to V_i$ के
        आव्यूहों के रूप में विघटित कीजिए और विभाज्यता शृंखला का
        उपयोग कीजिए।)}
  \item $\dim \mathcal C(u) \geq n$ निष्कर्ष निकालिए, जिसमें समता तभी और केवल तभी होती
        है जब $u$ चक्रीय हो ($s = 1$), और $u = \lambda\,\mathrm{id}$ के लिए $\dim\mathcal C(u)$
        परिकलित कीजिए: सूत्र के दोनों छोर।
  \item $\operatorname{diag}(J_2(0), J_1(0))$ के लिए क्रमविनिमेयी को $3\times3$ आव्यूहों के रूप में
        स्पष्ट परिकलित करके फ्रोबेनियस का सूत्र सीधे सत्यापित
        कीजिए।
\end{enumerate}

\medskip
\noindent\textbf{भाग III --- द्विक क्रमविनिमेयी प्रमेय।} मान
लीजिए $w \in \mathcal C(\mathcal C(u))$; हम $w \in K[u]$ सिद्ध करते हैं।
\begin{enumerate}[resume]
  \item दर्शाइए कि $K[u] \subseteq \mathcal C(\mathcal C(u))$, और यह कि प्रत्येक $w \in \mathcal C(\mathcal C(u))$ $u$ के साथ
        क्रमविनिमेय है --- अतः जो अंतर्वेशन सिद्ध करना है,
        $\mathcal C(\mathcal C(u)) \subseteq K[u]$, वह $w \in \mathcal C(u)$ का सच्चा परिष्कार है।
  \item पहले मान लीजिए $u$ \emph{चक्रीय} है, $V = K[u]x$। सीधे
        दर्शाइए कि $\mathcal C(u) = K[u]$ \emph{(क्रमविनिमेय करने वाले $v$ का
        $x$ पर मान लीजिए: किसी $P$ के लिए $v(x) = P(u)x$, और आधार
        $u^k x$ पर $v$ की $P(u)$ से तुलना कीजिए)}, और इस स्थिति
        में प्रमेय का निष्कर्ष निकालिए।
  \item अब व्यापक स्थिति पर लौटिए। प्रत्येक $i$ के लिए मान
        लीजिए $\pi_i\colon
        V \to V_i$ अन्य योज्यांशों के अनुदिश प्रक्षेप है।
        दर्शाइए कि $\pi_i \in \mathcal C(u)$, और इससे निष्कर्ष निकालिए कि $w$
        प्रत्येक $V_i$ को सुरक्षित रखता है और $u_i =
        u\restriction_{V_i}$ के साथ
        क्रमविनिमेय है; अब चक्रीय $u_i$ पर प्रश्न 11 लगाकर
        निष्कर्ष निकालिए: ऐसे बहुपद $Q_i$ हैं जिनके लिए $w\restriction_{V_i} = Q_i(u)\restriction_{V_i}$।
  \item अब $Q_i$ को एक ही बहुपद में जोड़ना शेष है। $i
        \leq j$ (अतः
        $P_i \mid P_j$) के लिए दर्शाइए कि $\eta_{ij} \colon
        V_j \to V_i$, $R(u)x_j \mapsto R(u)x_i$, एक सुपरिभाषित
        $K[X]$-समाकारिता है \emph{(जाँचने योग्य बात यह है कि
        $R(u)x_j = 0$ से $R(u)x_i = 0$ निकलता है)}, और यह कि $\tilde\eta_{ij} = \eta_{ij}\circ\pi_j$, जिसे अन्य
        योज्यांशों पर $0$ से विस्तारित किया गया हो, $\mathcal C(u)$ में
        है।
  \item $w\tilde\eta_{ij} = \tilde\eta_{ij}w$ का उपयोग करके $i \leq j$ के लिए $Q_i
        \equiv Q_j \pmod{P_i}$ दर्शाइए। इससे
        निष्कर्ष निकालिए कि $Q =
        Q_s$ सभी $i$ के लिए $Q \equiv Q_i \pmod {P_i}$ को
        संतुष्ट करता है, अतः प्रत्येक $V_i$ पर $w = Q(u)$, और इसलिए
        $V$ पर भी:
        \[
        \boxed{\ \mathcal C(\mathcal C(u)) = K[u].\ }
        \]
  \item (समापन) प्रमेय से निष्कर्ष निकालिए: यदि $v$ $u$ के
        साथ क्रमविनिमेय \emph{प्रत्येक} आव्यूह के साथ
        क्रमविनिमेय हो और $u$ चक्रीय हो, तो $v$ $u$ का एक
        बहुपद है; और ऐसा उदाहरण दीजिए जो दर्शाए कि $u =
        \mathrm{id}$, $n \geq 2$
        के लिए $\mathcal C(u) = K[u]$ विफल हो जाता है --- चक्रीयता ठीक कहाँ
        प्रवेश करती है?
\end{enumerate}

\medskip
\noindent\textbf{भाग IV --- समरूपता अपरिवर्तियों का लाभांश।}
परिमेय विहित रूप एक यंत्र है; यहाँ उसके पाँच शास्त्रीय उत्पाद
हैं।
\begin{enumerate}[resume]
  \item (परिवर्त) दर्शाइए कि प्रत्येक $M \in M_n(K)$ अपने परिवर्त ${}^tM$
        के समरूप है। \emph{($XI - M$ को स्मिथ रूप में लाने वाली
        संक्रियाएँ, परिवर्त लेने पर, $XI - {}^tM$ को उसी स्मिथ रूप में
        ले आती हैं: अपरिवर्ती समान।)}
  \item (समरूपता का अवरोहण) मान लीजिए $K \subseteq L$ एक क्षेत्र विस्तार
        है और $M, N \in M_n(K)$। दर्शाइए कि यदि $M$ और $N$ $L$ पर
        समरूप हों, तो वे $K$ पर भी समरूप हैं। \emph{($K[X]$ में
        परिकलित $XI - M$ का स्मिथ रूप $L[X]$ में भी स्मिथ रूप रहता है
        --- \omterm{thm:b3:modules:smith}{अपरिवर्ती गुणनखंड} क्यों नहीं बदलते?)} याद रखने योग्य
        परिणाम: $GL_n(\C)$ में संयुग्म दो वास्तविक आव्यूह $GL_n(\R)$ में
        भी संयुग्म होते हैं।
  \item (शून्यभावी वर्गीकरण) मान लीजिए $u$ शून्यभावी है।
        दर्शाइए कि शून्यभावी जॉर्डन खंडों में उसके विघटन में
        आकार $\geq k$ के खंडों की संख्या $\operatorname{rk}u^{k-1} - \operatorname{rk}u^k$ के बराबर है, और
        निष्कर्ष निकालिए: किसी भी क्षेत्र $K$ के लिए $M_n(K)$ के
        शून्यभावी वर्ग $n$ के विभाजनों के साथ एकैकी आच्छादक
        संगति में हैं। $M_5(K)$ में कितने शून्यभावी वर्ग हैं?
  \item (एक मूर्त युग्म) घात $< n$ के बहुपदों की समष्टि $K_{n-1}[X]$ पर
        क्रिया करने वाले अवकलन संकारक $D\colon P \mapsto P'$ के \omterm{thm:b3:modules:frobenius}{समरूपता अपरिवर्ती}
        ज्ञात कीजिए: (क) $K =
        \Q$ के लिए; (ख) $p < n$ के साथ $K = \mathbb F_p$ के
        लिए \emph{(अभिलक्षण $p$ में $(X^p)' = 0$: $\ker D^k$ परिकलित कीजिए
        और प्रश्न 18 का उपयोग कीजिए)}।
  \item ($GL_2(\mathbb F_q)$ के \omterm{ex:b3:groups:actions}{संयुग्मन वर्ग}) अपरिवर्ती गुणनखंडों का उपयोग
        करके दर्शाइए कि $GL_2(\mathbb F_q)$ का प्रत्येक वर्ग ठीक चार प्रकारों
        में से एक का है: केंद्रीय $aI$; $\mathbb F_q^\times$ में भिन्न-भिन्न
        अभिलक्षणिक मानों $a \neq b$ वाला विकर्णनीय; न्यूनतम बहुपद
        $(X - a)^2$ वाला अनर्धसरल; और \omterm{def:b3:rings:divisibility}{अखंडनीय} अभिलक्षणिक बहुपद वाला
        चक्रीय।
  \item प्रत्येक प्रकार के वर्ग गिनिए और निष्कर्ष निकालिए:
        $GL_2(\mathbb F_q)$ के ठीक $q^2 - 1$ \omterm{ex:b3:groups:actions}{संयुग्मन वर्ग} हैं। \emph{($\mathbb F_q$ पर
        मोनिक \omterm{def:b3:rings:divisibility}{अखंडनीय} द्विघात गिनिए; अक्रमित युग्म $\{a, b\}$;
        और स्मरण रखिए कि व्युत्क्रमणीयता अचर पदों को सीमित करती
        है।)}
  \item (चक्रीय होना प्ररूपी है) दर्शाइए कि $M \in M_2(\mathbb F_q)$ चक्रीय होने
        में विफल होता है तभी और केवल तभी जब $M$ अदिश हो, और
        इससे निष्कर्ष निकालिए कि $\mathbb F_q$ पर एकसमान यादृच्छिक
        $2\times2$ आव्यूह $1 - q^{-3}$ प्रायिकता से चक्रीय है। $M_n$ और बड़े
        $q$ के लिए समरूप अनुमान का कथन दीजिए (उपपत्ति अपेक्षित
        नहीं): अचक्रीय आव्यूह विरल होते हैं --- यही कारण है कि
        \cref{pb:b3:modules:1} के भाग III में केवल प्ररूपी स्थिति के आगे ही
        वास्तविक परिश्रम करना पड़ा।
\end{enumerate}

\medskip
\noindent\textbf{भाग V --- पूरक।}
\begin{enumerate}[resume]
  \item (क्रमविनिमेयी का \omterm{ex:b3:groups:actions}{केंद्र}) दर्शाइए कि बीजगणित $\mathcal C(u)$ का
        \omterm{ex:b3:groups:actions}{केंद्र} ठीक $K[u]$ है \emph{(भाग III के दोनों अंतर्वेशनों
        को मिलाइए)}। इससे निष्कर्ष निकालिए कि $\mathcal C(u)$ क्रमविनिमेय
        है तभी और केवल तभी जब $u$ चक्रीय हो --- इस प्रकार
        प्रश्न 8 की समता स्थिति विशुद्ध संरचनात्मक मार्ग से पुनः
        प्राप्त हो जाती है।
  \item (कौन-सी विमाएँ आती हैं?) फ्रोबेनियस के सूत्र से निष्कर्ष
        निकालिए कि प्रत्येक $u$ के लिए $\dim\mathcal C(u) \equiv n \pmod 2$। फिर $\dim V = 4$ वाले
        $\mathcal L(V)$ में $u$ के विचरण पर $\dim \mathcal C(u)$ द्वारा लिए गए मानों का
        ठीक-ठीक समुच्चय निर्धारित कीजिए: दर्शाइए कि वह $\{4, 6,
        8, 10, 16\}$ है
        \emph{($4$ में योग देने वाले घात अनुक्रम $n_1
        \leq \dots \leq n_s$ गिनिए
        और प्रत्येक को किसी शून्यभावी आव्यूह से साकार कीजिए)}।
        विशेष रूप से $12$ और $14$, यद्यपि सही सम-विषमता के
        हैं, प्राप्त नहीं होते: सम-विषमता की बाधा आवश्यक है पर
        पर्याप्त नहीं।
  \item ($GL_2(\mathbb F_3)$ का \omterm{cor:b3:groups:classeq}{वर्ग समीकरण}) $q = 3$ के लिए प्रश्न 20 के
        प्रत्येक \omterm{ex:b3:groups:actions}{संयुग्मन वर्ग} का आकार कक्षा--स्थायीकारी से
        परिकलित कीजिए: $GL_2(\mathbb F_q)$ में किसी \emph{चक्रीय} $M$ का
        \omterm{ex:b3:groups:actions}{केंद्रक} $K[M]$ का इकाई समूह है (प्रश्न 11)। तीनों
        अकेंद्रीय प्रकारों में $K[M]$ की पहचान कीजिए, $\mathbb F_3$ पर
        तीनों मोनिक \omterm{def:b3:rings:divisibility}{अखंडनीय} द्विघात सूचीबद्ध कीजिए, और \omterm{cor:b3:groups:classeq}{वर्ग
        समीकरण}
        \[
        48 = \abs{GL_2(\mathbb F_3)}
        = 2\cdot1 + 1\cdot12 + 2\cdot8 + 3\cdot6,
        \]
        को $2 + 1 + 2 + 3 = 8 = q^2 - 1$ वर्गों के साथ सत्यापित कीजिए, जैसा प्रश्न 21 ने
        बताया था।
\end{enumerate}
\end{problem}
